\documentclass[10pt]{article}
\usepackage{colm2026}
\usepackage{amsmath,amssymb}
\usepackage{booktabs}
\usepackage{graphicx}
\usepackage{multirow}
\usepackage{xcolor}
\usepackage{microtype}
\usepackage{natbib}
\usepackage{caption}
\usepackage{subcaption}

\title{Delphi-Psychometric Evaluation of 3D Generative Models:\\
       Combining Structured LLM Elicitation with\\
       Item Response Theory}

\author{Suki Wang \\ Roblox \\ \texttt{siqiwang@roblox.com}}

\date{Stanford AIMS Workshop on AI Measurement Science @ COLM 2026}

\begin{document}
\maketitle

\begin{abstract}
Benchmarks for 3D generative models aggregate automated metrics without principled
analysis of reliability, item difficulty, or judge consistency.
We introduce a framework combining three previously disconnected methodologies:
\textbf{Delphi-structured multi-round LLM judging} (for scalable, calibrated elicitation),
\textbf{Item Response Theory} (IRT) and \textbf{Many-Facet Rasch Modelling} (MFRM)
(for psychometric analysis of the resulting data), applied to Hi3DBench, a
public benchmark spanning six image-to-3D generation models and 510 prompts
across five quality criteria.
We report five empirical contributions:
(E1) the five object-level criteria load onto a single latent quality factor
     (1-factor CFA: CFI=0.95, RMSEA=0.09), validating within-level
     aggregation; a cross-level analysis adding Hi3DEval's four
     material-scoring dimensions confirms they form a psychometrically
     distinct factor (2F CFI=0.883 vs.\ 1F CFI=0.641, $\Delta$BIC$=-1{,}763$),
     so object-level and material-level scores should not be combined;
(E2) evaluation criterion, treated as the fixed benchmark-specified factor it
     is, accounts for 82\% of score variance ($\eta^2$; a random-effects
     estimate agrees at 84.7\%)—dwarfing model (1.4\%) and prompt (0.9\%),
     with previously-omitted model$\times$criterion and prompt$\times$criterion
     interactions comparably small (2.1\%, 1.5\%); per-criterion standardization
     directly confirms this is a scale artefact of the automated metrics,
     collapsing criterion's share to $\approx$0\% and lifting model's share
     nearly 10-fold. Delphi LLM judging on a uniform 0--9 scale cross-validates
     this finding by making model quality the dominant source (37\%), and
     achieves strong reliability (ICC$=0.81$--$0.88$; inter-judge SD reduced
     68\% across Delphi rounds);
(E3) in a flipped 2PL IRT orientation (prompts as persons, models as items,
     $N=470$ eligible prompts), model stringency spans a $>8$-logit range
     (HunyuanDiT-3D $b=-5.29$ to CRM $b=2.83$); prompt generatability
     $\hat{\theta}_p\in[-1.07,1.14]$ with 75 near-ceiling prompts ($p\geq0.8$);
(E4) descriptive sensitivity analysis reveals 71.6\% of items show
     $|\text{gap}|\geq0.40$ in pass rates between TripoSR and the five
     Objaverse-trained models; because Group~B contains exactly one model,
     Objaverse exposure and TripoSR-specific factors are perfectly confounded
     and cannot be separated;
(E5) an IRT-$\theta$-stratified 30-item \emph{tinyEval3D} subset recovers the
     full-benchmark model ranking exactly (Spearman $\rho=1.0$, LOO-CV
     $\rho=0.954$), reducing evaluation cost by 94\%, with model rankings
     confirmed by Delphi LLM judging ($\rho=0.90$ agreement with automated
     metric rankings).
Our framework exposes structural inefficiencies in current 3D benchmarks and
provides a reusable methodology for constructing psychometrically grounded
evaluations.
\end{abstract}

%% ─────────────────────────────────────────────────────────────────────────────
\section{Introduction}
\label{sec:intro}

The rapid proliferation of text-to-3D and image-to-3D generative models has
produced an equally rapid proliferation of evaluation benchmarks
\citep{eval3d2025,hi3deval2025,t3bench2023,gen3deval2024}.
These benchmarks share a common methodology: aggregate automated metric scores
into weighted sums, report raw percentage-agreement with human judgment, and
treat all evaluation items as equally informative.
This methodology has well-known limitations in adjacent fields.
Sixty years of psychometric research in educational testing have established
that: (i) score aggregation is only valid when items measure a unidimensional
construct, (ii) items differ systematically in difficulty and discriminability,
(iii) rater effects must be separated from item effects, and (iv) a small set
of maximally informative items can often substitute for large benchmarks.
\textbf{Item Response Theory} (IRT) formalizes (ii) by modeling the
probability of a correct/passing response as a function of a latent person
ability and item-specific difficulty and discrimination parameters, rather
than treating every item as an equally-weighted point
\citep{embretson2000}. \textbf{Many-Facet Rasch Measurement} (MFRM)
extends this to designs with additional measurement facets beyond
persons and items---most relevantly, raters---so that a rater's severity
is estimated and separated from the quantity being measured rather than
folded into it \citep{linacre1994,eckes2011}. We apply both: IRT (via a
graded-response and a 2PL model) to model quality items, and an
MFRM-style crossed random/fixed-effects decomposition to separate model,
prompt, criterion, and judge facets.

Parallel work in text-LLM evaluation has begun applying psychometric methods
\citep{tinybenchmarks2024}.
However, this work has not reached 3D generative evaluation, where the
multi-modal, multi-view nature of the task and the absence of ground-truth
labels make reliability analysis especially critical.

At the same time, LLM-as-judge approaches have become standard in generative AI
evaluation \citep{llmjudge2023}, but are typically applied as
one-shot prompts without any structured disagreement resolution.
When the same judge is queried repeatedly, or when judges are shown each
other's outputs with identity visible, they are susceptible to
\emph{authority bias}---conforming to a higher-status model's assessment---and
\emph{anchor conformity}---converging toward a revealed group median without
genuinely re-examining the stimulus.
The Delphi method \citep{dalkey1963delphi}, developed at RAND in the 1960s for
policy forecasting, provides a protocol for iterative expert consensus that
directly addresses both failure modes and has never been applied to LLM
evaluation.

This paper contributes the \textbf{Delphi-Psychometric triangle}:
Delphi structures the elicitation; LLMs scale it; psychometrics analyses the
resulting data.
We demonstrate the framework on Hi3DBench \citep{hi3deval2025}, producing five
empirical findings (E1–E5) that reveal structural properties of the benchmark
invisible to standard analysis.

%% ─────────────────────────────────────────────────────────────────────────────
\section{Related Work}
\label{sec:related}

\paragraph{3D generative evaluation.}
Eval3D \citep{eval3d2025} introduces five automated metrics covering geometric,
structural, semantic, aesthetic, and text-alignment dimensions.
Hi3DEval \citep{hi3deval2025} extends this with hierarchical human-aligned
scoring across object- and part-level properties.
T3Bench \citep{t3bench2023} and Gen3DEval \citep{gen3deval2024} provide
complementary text-to-3D benchmarks.
All use raw metric averages; none apply psychometric analysis.

\paragraph{Psychometrics for AI evaluation.}
tinyBenchmarks \citep{tinybenchmarks2024} applies Item Response Theory to
text-LLM benchmarks, showing that $\sim$100 items can recover full-benchmark
rankings.
None of these works addresses 3D evaluation.

\paragraph{LLM-as-judge and Delphi.}
LLM-as-judge / MT-Bench \citep{llmjudge2023} demonstrates
LLM raters as scalable alternatives to human evaluation.
The Delphi method \citep{dalkey1963delphi,rowe1999delphi} is an established
protocol for iterative expert consensus in high-uncertainty domains.
To our knowledge, Delphi has not previously been applied to LLM judging.

%% ─────────────────────────────────────────────────────────────────────────────
\section{Methods}
\label{sec:methods}

\subsection{Dataset: Hi3DBench}
We use Hi3DBench \citep{hi3deval2025}, a public benchmark comprising rendered
images of 3D assets generated by six models: CRM, HunyuanDiT-3D, SPARD,
TRELLIS, TripoSR, and Unique3D.
We use 510 prompts from the shared prompt library, yielding $510 \times 6 = 3{,}060$
(prompt, model) pairs.
Each pair is evaluated on five quality criteria, aligned to Eval3D dimensions:
\textbf{geometric consistency}, \textbf{structural consistency},
\textbf{semantic consistency}, \textbf{aesthetics}, and
\textbf{text-3D alignment}.
Psychometric analyses (E1--E5) use the Hi3DBench automated scores
(geometry\_score, geo\_detail\_score, geo\_texture\_score, texture\_score,
alignment\_score) mapped to Eval3D criteria; these were validated against
human judgments in the original paper (Table~1, \citealt{hi3deval2025}).
For Delphi LLM calibration, we sample 100 prompts from the same library and
apply the protocol to the five models with complete rendered archives
(HunyuanDiT-3D, SPARD, TRELLIS, TripoSR, Unique3D), yielding $500$
(prompt, model) pairs rated by three LLM judges across up to three rounds.

\subsection{Delphi-Structured LLM Judging Protocol}
\label{sec:delphi}

To collect LLM-based ratings, we implement a three-round Delphi protocol
with three vision-capable LLM judges: \textbf{Claude Opus~4.6}
(Anthropic), \textbf{GPT-4.1} (OpenAI), and \textbf{Gemini~2.5~Pro} (Google).

\paragraph{Round structure.}
In Round~1 (independent rating), each judge receives a multi-view rendering
of the asset along with the text prompt, criteria definitions, and
VQA questions for the text-3D alignment dimension, and returns a structured
JSON rating (0–9 integer scale per criterion) with chain-of-thought reasoning.

In Round~2 (informed re-rating), each judge sees its own Round~1 scores, the
group median and IQR across judges, an anonymized dissent summary
generated by a fourth Claude~Opus~4.7 call, and the Hi3DBench automated
score for the same asset (framed as a baseline to consider, not a target
to match; judges are asked to explain rather than revise when their rating
diverges by more than 2 points).
Judges may revise or confirm each score; agreement with the automated score
is never required or scored.

Round~3 follows the same structure as Round~2 with Round~2 statistics.
We apply the \citet{delphi_convergence} stopping rule: if inter-judge SD
drops below 1.0 on $\geq 4$ of 5 criteria after Round~2, Round~3 is skipped
to reduce cost.

\paragraph{Dissent summary.}
Between rounds, we summarize inter-judge reasoning disagreements using a
dedicated LLM call (Claude Opus~4.7) with an anonymizing prompt, producing a
3--4 sentence summary that preserves disagreement content while removing
judge identity.

\paragraph{Anti-sycophancy design.}
Three design choices specifically address sycophancy failure modes present in
standard LLM-as-judge.
First, \textbf{independent Round~1}: each judge scores the asset before seeing
any other judge's output, anchoring their position to the stimulus rather than
to a social prior.
Second, \textbf{distinct evaluative personas}: Claude acts as a conservative
technical artist focused on surface quality and geometry, GPT as an analytical
CV researcher focused on multi-view consistency, and Gemini as a UX designer
focused on semantic clarity and user perception.
These personas create structurally different failure modes, so genuine
disagreements reflect substantive differences in evaluative framework rather
than random noise.
Third, \textbf{anonymized dissent}: judge identity is stripped from the
inter-round summary, preventing judges from simply deferring to the model they
perceive as most authoritative.
Residual anchoring to the group median---visible in \texttt{group\_stats\_json}---cannot be eliminated by design, but requiring written justification for any revision creates an epistemic check on pure conformity.
Exposing the automated score from Round~2 onward is a second, distinct
anchoring channel: because every judge sees the same external reference
point, part of the inter-judge convergence we report (\S\ref{sec:e2}) may
reflect shared regression toward that anchor rather than purely
independent deliberation, and the agreement we later report between
Delphi and automated rankings is accordingly not fully independent
confirmation. We retain this design because removing the anchor would
require re-running the full protocol, and because the Delphi stream still
disagrees with the automated stream on the top-ranked model
(\S\ref{sec:discussion}) despite the anchor pushing toward agreement---if
anything, that disagreement is stronger evidence of a genuine measurement
difference than it would be under a fully blind protocol.

\subsection{Psychometric Analysis Pipeline}
All analyses use R~4.5 with packages \texttt{mirt}, \texttt{lme4},
\texttt{lavaan}, \texttt{psych}, \texttt{difR}, and \texttt{irr}.
Code and data will be released upon publication.

\paragraph{E1 — Factor structure (Multidimensional IRT).}
We apply Horn's parallel analysis, exploratory factor analysis (EFA) with
oblimin rotation, confirmatory factor analysis (CFA), and graded-response
MIRT to the (model $\times$ prompt) $\times$ 5-criterion score matrix,
assessing whether the five criteria measure a single latent quality dimension.

\paragraph{E2 — Variance decomposition (MFRM).}
We fit stream-specific mixed models. Criterion is a fixed, benchmark-specified
set of five rubric dimensions rather than a sample from a population of
criteria, so we treat it as a \emph{fixed} facet ($\beta_c$); model, prompt,
judge, and round remain random facets, consistent with generalizability-theory
practice for designs with a mix of fixed and random facets
\citep[e.g.,][]{cronbach1972,brennan2001}. Interactions of the fixed criterion
factor with a random facet are themselves random effects.
For the \textbf{automated-score stream} (no judges or rounds):
\[
y_{mpc} = \mu + \beta_c + u_m + u_p + u_{mp} + u_{mc} + u_{pc} + \varepsilon
\]
For the \textbf{Delphi stream} (three judges, up to three rounds),
two additional facets enter:
\[
y_{mpjcr} = \mu + \beta_c + u_m + u_p + u_j + u_r + u_{mp} + u_{mc} + u_{pc} + \varepsilon
\]
In both equations $m$=model, $p$=prompt, $c$=criterion, $u_{mp}$ is a
model$\times$prompt interaction, and $u_{mc}$, $u_{pc}$ are
model$\times$criterion and prompt$\times$criterion interactions (previously
omitted); $j$=judge and $r$=round apply to the Delphi stream only.
The fixed criterion effect's share of total variance is reported as $\eta^2$
from a one-way ANOVA on criterion (unambiguous here since criterion is the
only fixed factor); remaining facets are REML variance components from the
mixed model, rescaled to sum to $1-\eta^2$ so the full decomposition adds to
100\%. As a robustness check, we also report the alternative in which
criterion is instead treated as random (matching a common but, for a fixed
rubric, less appropriate convention), and a version in which each criterion's
scores are $z$-scored (mean 0, SD 1) before pooling, to isolate the
contribution of raw scale differences across criteria from genuine model or
prompt signal.
We compute variance components as percentages of total variance, estimate
the G-coefficient (generalizability), and quantify the Delphi convergence
effect as the proportional reduction in judge-level variance from Round~1
to Round~3.
Inter-rater reliability is reported as ICC(A,$k$)---two-way random-effects,
absolute agreement, average-of-$k$-raters, in the McGraw--Wong
\citep{mcgraw1996} naming (equivalently ICC(2,$k$) in
Shrout--Fleiss~\citeyearpar{shrout1979} notation)---the appropriate form
given all three judges rate every item (a fully crossed design) and we
report the reliability of the mean of three judges' ratings, not of a
single judge.

\paragraph{E3 — Model stringency and prompt generatability (flipped 2PL IRT).}
We dichotomize \textbf{geometric consistency} at the midpoint (threshold=5,
sensitivity at 3,4,6,7) to produce a binary (model $\times$ prompt) matrix;
this criterion alone spans a wide enough native range to support a
midpoint split, unlike the other four criteria, whose automated-scoring
ranges are capped well below the midpoint (\S\ref{sec:e2}) and would make
any fixed threshold degenerate. CRM's automated scores are unavailable for
263 of the 510 prompts (generation failures); these are scored as automatic
fails on this criterion rather than excluded, so that a model's inability
to generate a scoreable asset for a prompt counts against it as a real
benchmark outcome instead of silently shrinking its eligible sample.
Because IRT stability requires $N\geq200$ persons and we have only 6 models,
we \emph{transpose} the matrix—placing 470 eligible prompts as persons and
6 models as items—and fit a 2PL IRT model.
This orientation yields stable model parameters ($b_m$: stringency;
$a_m$: discrimination power) and prompt ability estimates
$\hat{\theta}_p$ (generatability), computed via EAP scoring.

\paragraph{E4 — Training-Data Sensitivity Analysis.}
We group models by Objaverse training exposure
(Group~A: five Objaverse-trained models; Group~B: one non-Objaverse model, TripoSR).
Because standard inferential DIF tests (Mantel–Haenszel, logistic regression
\citealt{difR2010}) require $n \geq 10$ per group, we conduct a
\emph{descriptive sensitivity analysis}: for each item, we compute the gap in
pass rates between the two groups ($\Delta p = p_\text{Objaverse} - p_\text{non-Objaverse}$)
and report the distribution of $|\Delta p|$ as a measure of training-data sensitivity.

\paragraph{E5 — Adaptive 30-item subset (tinyEval3D).}
We compare four 30-item selection strategies: (i) IRT-$\theta$-stratified
sampling (equal items from each quintile of the estimated prompt generatability
$\hat{\theta}_p$ from the flipped 2PL), (ii) random sampling (1{,}000 draws),
(iii) classical difficulty-stratified sampling (quintiles of $p$-value),
and (iv) top-$N$ by point-biserial discrimination.
We report Spearman $\rho$ between subset-derived and full-benchmark model
rankings, validated by leave-one-model-out cross-validation.

%% ─────────────────────────────────────────────────────────────────────────────
\section{Results}
\label{sec:results}

\subsection{E1: Factor Structure}
\label{sec:e1}

CFA fit statistics (Table~\ref{tab:cfa}) show that a 1-factor model achieves
CFI=0.95, RMSEA=0.087 (SRMR=0.037), consistent with a primarily
unidimensional latent quality construct; the 1-factor model also wins on CFA BIC
($\Delta$BIC$=+6.7$ for 2F vs.\ 1F).
EFA and graded-response MIRT both favour a 2-factor solution
(EFA BIC$=-3.87$ for 2F vs.\ $+71.6$ for 1F; MIRT BIC$=15{,}577$ for $k=2$
vs.\ $15{,}594$ for $k=1$), where Factor~1 captures structural coherence
(structural\_consistency loading~0.86) and Factor~2 captures content accuracy
(text\_3d\_alignment~0.64, geometric\_consistency~0.52, semantic\_consistency~0.44).
Together, these results indicate that while a single global quality score is
a reasonable approximation (CFA 1F: CFI=0.95), the criteria have a meaningful
two-factor substructure.

\begin{table}[h]
\centering
\small
\caption{CFA Fit Indices. 1F = one-factor model (global quality); 2F = geometry vs.\
semantic-aesthetic split. RMSEA $<$ 0.08 indicates acceptable fit; CFI $>$ 0.95 is good.}
\label{tab:cfa}
\begin{tabular}{lrrrrl}
\toprule
Model & CFI & TLI & RMSEA & SRMR & $\Delta$BIC (vs.\ 1F) \\
\midrule
1F & 0.953 & 0.905 & 0.087 & 0.037 & — \\
2F & 0.953 & 0.882 & 0.097 & 0.036 & +6.7 \\
\bottomrule
\end{tabular}
\end{table}

\textbf{Cross-level extension: object-scoring vs.\ material-scoring.}
Hi3DEval includes a second evaluation layer---material scoring---covering
four physical-property dimensions (Details \& Complexity, Colorfulness \&
Saturation, Consistency \& Artifacts, Material Plausibility; scale 0--4).
We extend the factor analysis to all nine criteria jointly on the 1{,}730
(model, prompt) pairs for which both layers are available (4 models: CRM,
HunyuanDiT-3D, SPARD, TripoSR).

Table~\ref{tab:cfa_cross} shows that a single-factor model fails badly across
the nine criteria (CFI=0.641, RMSEA=0.237), while a two-factor model
separating the object-level and material-level blocks fits substantially
better (CFI=0.883, RMSEA=0.138, $\Delta$BIC$=-1{,}763$ vs.\ 1F).
EFA $k=2$ produces textbook-clean loadings: all four material criteria load
exclusively on Factor~1 (loadings 0.67--0.86, cross-loadings $<0.10$) and
all five object criteria load exclusively on Factor~2 (loadings 0.55--0.75,
cross-loadings $<0.10$).
This separation is further confirmed by correlation structure: mean
within-block $|r|$ is 0.47 (object) and 0.62 (material), versus only 0.28
between blocks.

\begin{table}[h]
\centering
\small
\caption{CFA fit for the 9-criterion (5 object $+$ 4 material) joint model
on $N=1{,}730$ observations. The hierarchical 2F model separates
object-level and material-level dimensions; $\Delta$BIC relative to 1F.}
\label{tab:cfa_cross}
\begin{tabular}{lrrrr}
\toprule
Model & CFI & RMSEA & SRMR & $\Delta$BIC \\
\midrule
1F (global quality)               & 0.641 & 0.237 & 0.141 & — \\
2F (object $|$ material) & 0.883 & 0.138 & 0.054 & $-1{,}763$ \\
\bottomrule
\end{tabular}
\end{table}

\textbf{Implication.}
Within each evaluation level the five object-level criteria are approximately
unidimensional (CFA 1F: CFI=0.95, as reported above), validating
weighted-sum aggregation \emph{within} a level.
Across levels, however, object and material dimensions are psychometrically
distinct constructs that load onto separate factors with near-zero
cross-loadings---confirming Hi3DEval's hierarchical validity claim
empirically.
A single aggregate score combining both layers would conflate two
genuinely different aspects of 3D quality and should be avoided.
Criterion-level reporting remains valuable for diagnosing failure modes;
the two-factor substructure within the object layer (structural coherence
vs.\ content accuracy) provides additional diagnostic resolution beyond the
object-level aggregate.
Note that the 3{,}060 (prompt, model) observations are nested within 510
prompts and 6 models; a full multilevel CFA would be needed to formally
account for this clustering, which is left for future work.

\subsection{E2: Variance Decomposition}
\label{sec:e2}

The mixed-effects MFRM with criterion as a fixed factor (Table~\ref{tab:vc})
shows that \textbf{evaluation criterion} accounts for 82.0\% of total score
variance ($\eta^2$)---dwarfing model (1.4\%), prompt (0.9\%), and the
model$\times$prompt interaction (2.0\%); the two previously-omitted
interactions of criterion with the random facets, model$\times$criterion and
prompt$\times$criterion, are each comparably small (2.1\%, 1.5\%), so their
exclusion in a simpler model would not have changed the qualitative
conclusion, though we report them for completeness.
Treating criterion instead as a random facet (the more common but, for a
fixed 5-item rubric, less appropriate convention) gives a closely agreeing
estimate of 84.7\% (supplementary material), indicating the finding is
not an artefact of that modeling choice.
Residual variance (10.1\%) captures unexplained item-level noise.
This criterion-dominance is itself a scale artefact of the automated metrics:
each criterion uses a different scoring function with its own mean and range,
so \emph{which} criterion is measured swamps \emph{which} model produced the
asset. We confirm this directly, rather than only by contrast with the Delphi
stream below: standardizing each criterion to mean~0/SD~1 before pooling
collapses criterion's share to $\approx$0\% by construction and reallocates
the removed variance overwhelmingly to prompt (20.7\%) and model (14.8\%)---a
roughly 10-fold increase in model's share once the scale artefact is removed
(full table in supplementary material).
The 82.0\% figure should therefore be read as an illustration of artefact
severity rather than a psychometrically meaningful partition of quality
variance.

\begin{table}[h]
\centering
\small
\caption{Variance components from the Many-Facet Rasch Model (Hi3DBench
automated scores), with criterion treated as a fixed effect
(\S\ref{sec:e2}). Criterion-level scale artefacts dominate (82.0\%); genuine
model and prompt effects, including their interactions with criterion, are
small in comparison. All rows sum to 100\%; no residual double-counting.}
\label{tab:vc}
\begin{tabular}{lr}
\toprule
Facet & \% of Total \\
\midrule
Criterion (fixed effect)         & 82.0\% \\
Residual                         & 10.1\% \\
Model$\times$Criterion           &  2.1\% \\
Model$\times$Prompt              &  2.0\% \\
Prompt$\times$Criterion          &  1.5\% \\
Model                            &  1.4\% \\
Prompt                           &  0.9\% \\
\midrule
Total                            & 100.0\% \\
\bottomrule
\end{tabular}
\end{table}

\textbf{Delphi cross-validation.}
To verify that the criterion-dominance finding is an artefact rather than a
genuine psychometric property, we ran the same fixed-criterion MFRM on the
Delphi LLM ratings, which are collected on a uniform 0--9 scale for all five
criteria.
When scale artefacts are removed, \textbf{model quality} becomes the dominant
source of variance (37.3\%), with criterion dropping to 6.0\% ($\eta^2$;
6.6\% under the random-effects estimate) and judge to 1.7\%
(supplementary material).
Unlike in the automated stream, prompt$\times$criterion is not negligible
here (5.4\%, larger than criterion's own main effect and than
model$\times$criterion at 1.3\%), suggesting some prompts are differentially
easy or hard depending on which criterion is being judged (e.g.\ geometrically
simple but semantically ambiguous); prompt, round, model$\times$prompt, and
residual together account for the remaining $\approx$49\%.
Per-criterion standardization changes this picture only modestly (model
37.3\%$\to$41.2\%)---in sharp contrast to the automated stream's near
10-fold shift---because the Delphi scale is already uniform across criteria
by design, confirming that the scale artefact is specific to the automated
scoring functions rather than an inherent property of the underlying quality
construct.
Delphi judge reliability is high (Krippendorff $\alpha=0.56$--$0.70$ per
criterion; ICC$=0.81$--$0.88$), and inter-judge SD drops by 67.8\% from
Round~1 to Round~3 (1.663 $\to$ 0.535), validating the convergence of
the elicitation protocol.
The G-coefficient for model rankings is 0.986, indicating excellent
generalizability of automated benchmark scores across the prompt universe.

\subsection{E3: Model Stringency and Prompt Generatability}
\label{sec:e3}

Of 510 prompts, 40 have zero variance (all six models pass or all fail on
geometric consistency at threshold=5) and are excluded, leaving $N_p=470$
eligible prompts.
Treating these prompts as IRT ``persons'' and the six models as ``items''
satisfies the $N\geq200$ stability requirement (the classical threshold
for reliable 2PL calibration).

The 2PL model converges cleanly and reveals a striking spread in model
stringency (Table~\ref{tab:irt}):
HunyuanDiT-3D is the easiest model to satisfy ($b=-5.29$; 87.5\% prompt
pass rate), while CRM is the hardest ($b=2.83$; 11.0\% pass rate)---an
$>8$-logit range.
TRELLIS ($b=-0.99$) and SPARD ($b=-0.34$) cluster near the midpoint,
with TripoSR ($b=0.65$) and Unique3D ($b=1.45$) in the upper-difficulty
half.
Discrimination $a_m$ ranges from 0.48 (HunyuanDiT-3D, near-ceiling pass
rate reduces discriminability) to 1.08 (SPARD), indicating that SPARD and
TRELLIS are the most informative models for separating easy from hard prompts.

Prompt generatability $\hat{\theta}_p$ spans $[-1.07,\ 1.14]$, confirming a
wide difficulty range; 75 prompts (16\%) are near-ceiling ($p\geq0.8$),
meaning the benchmark does contain a meaningful tail of easy items.

\begin{table}[h]
\centering
\small
\caption{Flipped 2PL IRT model parameters. $b_m$ = model stringency
(lower = easier; more prompts pass). $a_m$ = discrimination (how sharply
model $m$ separates easy from hard prompts). Pass rate = proportion of
prompts for which model $m$ scores $\geq5$ on geometric consistency.}
\label{tab:irt}
\begin{tabular}{lrrl}
\toprule
Model & $b_m$ (stringency) & $a_m$ (discrimination) & Pass rate \\
\midrule
HunyuanDiT-3D & $-5.29$ & 0.48 & 87.5\% \\
TRELLIS       & $-0.99$ & 0.95 & 66.3\% \\
SPARD         & $-0.34$ & 1.08 & 55.7\% \\
TripoSR       & $+0.65$ & 0.57 & 41.0\% \\
Unique3D      & $+1.45$ & 0.63 & 30.6\% \\
CRM           & $+2.83$ & 0.94 & 11.0\% \\
\bottomrule
\end{tabular}
\end{table}

\textbf{Implication.}
The flipped orientation solves the small-$N$ instability of standard IRT
(6 models as persons would leave each item with only 6 binary observations)
by exploiting the much larger prompt pool.
The resulting prompt-generatability scores $\hat{\theta}_p$ directly inform
the IRT-$\theta$-stratified item selection in E5.

\subsection{E4: Training-Data Sensitivity}
\label{sec:e4}

Because only one model (TripoSR) lacks Objaverse training, inferential DIF
tests are underpowered and inappropriate here (minimum $n\geq10$ per group
required).
We instead report a descriptive sensitivity analysis comparing per-item pass
rates between the five Objaverse-trained models and TripoSR.

Table~\ref{tab:dif} summarises the distribution of $|\Delta p|$.
\textbf{71.6\%} of items show $|\Delta p|\geq0.40$, and \textbf{14 items} (2.7\%)
show complete reversals ($|\Delta p|=1.0$, i.e., all Objaverse models pass
while TripoSR fails or vice versa).

\begin{table}[h]
\centering
\small
\caption{Training-data sensitivity: distribution of $|\Delta p|$ (pass-rate gap between
Objaverse-trained and non-Objaverse model groups) across 510 items.}
\label{tab:dif}
\begin{tabular}{lrrl}
\toprule
$|\Delta p|$ range & Items & \% & Sensitivity \\
\midrule
$[0.0,\ 0.2)$  &  40 &  7.8\% & Negligible \\
$[0.2,\ 0.4)$  & 105 & 20.6\% & Mild \\
$[0.4,\ 0.6)$  & 177 & 34.7\% & Moderate \\
$[0.6,\ 0.8)$  & 118 & 23.1\% & Large \\
$[0.8,\ 1.0]$  &  70 & 13.7\% & Severe \\
\bottomrule
\end{tabular}
\end{table}

\textbf{Implication.}
The $|\Delta p|$ distribution is better described as \emph{TripoSR sensitivity}
than as \emph{Objaverse sensitivity}, because Group~B consists of exactly one
model.
With $N_\text{non-Objaverse}=1$, the ``training-data effect'' and
``TripoSR-specific effect'' (architecture, pre-training corpus, fine-tuning
choices) are \emph{perfectly confounded}: there is no mathematical way to
separate them.
Any gap we observe could reflect (a) a genuine Objaverse-vs-non-Objaverse
difference, (b) something idiosyncratic to TripoSR's design that has nothing
to do with Objaverse, or (c) both simultaneously.
We therefore concede that the analysis cannot be interpreted as evidence for
or against training-data bias in the benchmark; it is, at best, a flag that
TripoSR behaves differently from the other five models on 71.6\% of items.
Resolving the confound requires adding at least two or three more
non-Objaverse models (e.g.\ models trained on ShapeNet or G-Objaverse
subsets) to give Group~B sufficient size for formal Mantel--Haenszel DIF
testing.

\subsection{E5: tinyEval3D Adaptive Subset}
\label{sec:e5}

Table~\ref{tab:subset} compares four 30-item selection strategies.
Under leave-one-model-out cross-validation (LOO-CV)---the more conservative
estimate of practical reliability when $N=6$ models calibrate the item
selection---both the IRT-$\theta$-stratified and classical difficulty-stratified
strategies achieve Spearman $\rho=0.954$ (see supplementary material).
On the full benchmark (all 6 models used for calibration),
\textbf{IRT-$\theta$-stratified} selection achieves $\rho=1.0$
(exact ranking recovery) and \textbf{classical difficulty-stratified}
achieves $\rho=0.943$.
Notably, random sampling achieves a high mean $\rho=0.974\pm0.036$, reflecting
that with only 6 models any 30 prompts tend to preserve rankings;
top-discrimination selection ($\rho=0.878$) underperforms all other strategies.
This tempers how strongly E5 should be read: with model rankings this
stable, the honest claim is not that IRT-based selection is
\emph{necessary} to recover them---classical difficulty-stratification,
which needs no IRT model at all, achieves nearly identical recovery
($\rho=0.943$, identical LOO-CV$=0.954$)---but that \emph{some} form of
difficulty-aware stratification is what does the work, and IRT-$\theta$
provides that stratification without an ad hoc search over random draws
or dependence on how separated this particular set of 6 models happens to
be. We expect the gap between stratified and random selection to widen as
the model pool grows and rankings become less trivially separable.

\begin{table}[h]
\centering
\small
\caption{tinyEval3D: Ranking Recovery by 30-Item Selection Strategy.
Spearman $\rho$ between subset-derived and full-benchmark model rankings
(30 items from 510; $N=6$ models).}
\label{tab:subset}
\begin{tabular}{lrrr}
\toprule
Strategy & Spearman $\rho$ & SD ($\rho$) & LOO-CV $\rho$ \\
\midrule
IRT-$\theta$-stratified (\textbf{tinyEval3D}) & \textbf{1.000} & — & \textbf{0.954} \\
Classical difficulty-stratified & 0.943 & — & 0.954 \\
Random (mean over 1{,}000) & 0.974 & 0.036 & — \\
Top-discrimination & 0.878 & — & — \\
\bottomrule
\end{tabular}
\end{table}

The IRT-$\theta$-stratified strategy uses prompt generatability scores
$\hat{\theta}_p$ from the flipped 2PL (Section~\ref{sec:e3}) to ensure
the 30-item subset covers the full difficulty range in a
psychometrically principled way.
Classical difficulty-stratified selection achieves nearly equivalent LOO-CV
performance using only classical $p$-values, making it the preferred strategy
when IRT theta is unavailable.
A 30-item \textbf{tinyEval3D} subset reduces evaluation cost by 94\% while
preserving near-exact model ranking fidelity, enabling rapid iteration during
model development.

%% ─────────────────────────────────────────────────────────────────────────────
\section{Discussion}
\label{sec:discussion}

\paragraph{Criterion dominates; rater and prompt matter less; IRT agrees
with the naive rating.}
The 82\% criterion variance share (E2, fixed-effect $\eta^2$; 84.7\% under a
random-effects estimate) is the most striking finding in this paper: a
given prompt's score is determined far more by \emph{which} criterion is
being measured than by \emph{which} model generated it.
This is a scale artefact, not a property of quality itself: per-criterion
standardization on the automated data collapses criterion's share to
$\approx$0\% directly, and on Delphi's independently-scaled uniform 0--9
stream criterion drops to 6.0\% while model quality becomes the dominant
source (37.3\%); judge (1.7\%) and prompt (12.9\%) are smaller than model
but not negligible.
The implication for benchmark design is direct: raw aggregation of
criterion-specific automated metrics conflates genuine quality differences
with scale differences; evaluators should either normalise metrics per
criterion or use a common-scale judging approach (such as Delphi LLM
elicitation) before aggregating scores.
Separately, we asked whether accounting for prompt difficulty and
criterion informativeness through IRT actually changes which model looks
best, or is equivalent to a naive average. Extracting IRT ability
estimates ($\theta$, EAP) from E1's graded-response MIRT---whose criterion
discrimination parameters weight informative criteria more than saturated
ones, unlike a naive equal-weighted mean---reproduces the identical
ranking to the naive mean in \emph{both} streams, and this holds
separately for the structural-coherence and content-accuracy sub-factors
of E1's two-factor solution. Within a single stream, IRT reweighting does
not change who wins; a naive mean is just as good a summary. Where IRT
earns its keep is in adjudicating \emph{closer} comparisons across
multiple dimensions and, as we show next, across measurement instruments
entirely.

\paragraph{Delphi helps with convergence, not necessarily validity---but
that is exactly when it is useful.}
Standard LLM-as-judge elicits a single round of ratings with no structured
disagreement resolution, leaving it susceptible to authority bias and
anchor conformity when judges see each other's outputs. Our Delphi
protocol mitigates this via anonymized dissent summaries, persona
diversity, and mandatory justification for revisions, and achieves high
inter-judge reliability (ICC$=0.81$--$0.88$) with a 67.8\% SD reduction
across rounds (1.663 $\to$ 0.535). But convergence is not the same claim as
validity: falling disagreement could reflect genuine deliberation, or
could partly reflect judges regressing toward a shared anchor---the group
median, and, as we disclose in \S\ref{sec:delphi}, the automated score
itself, shown from Round~2 onward. We cannot fully separate these without
independent human ground truth, which we lack at item level (Hi3DBench's
automated scores were validated against humans by its original authors,
not by us). This is precisely the situation where a second, independent
measurement stream is most useful: not because Delphi is correct where
automated scores might be wrong, but because it is a differently-biased
instrument, and where two differently-biased instruments still mostly
agree (Spearman $\rho=0.90$ overall model-ranking agreement; item-level
$r=0.22$--$0.53$, weakest on structural consistency, where LLM judges are
least confident assessing fine-grained geometric coherence from rendered
images), that agreement is evidence neither is pure noise. Where they
disagree outright, as below, that disagreement is diagnostic in a way
neither instrument alone could reveal.

\paragraph{IRT to adjudicate disagreement: HunyuanDiT-3D or TRELLIS?}
On automated scores, HunyuanDiT-3D ranks first (mean composite 2.95 vs.\
TRELLIS' 2.74); on Delphi ratings, TRELLIS ranks first (4.93 vs.\
HunyuanDiT-3D's 4.64)---the two streams disagree on which model is best.
(The overall $\rho=0.90$ agreement figure is consistent with this: with
only 5 models rated by both streams, one adjacent rank swap at the top
already produces $\rho=0.90$ exactly.) Which should we believe?
First, we rule out a dimension-weighting artefact: perhaps each model wins
different criteria, and ``best'' depends on how the five criteria are
combined. It does not---the IRT-weighted ranking from the previous
paragraph is identical to the naive ranking in both streams, and each
winning model wins on \emph{every} axis (structural coherence and content
accuracy) within its own stream, not a subset. Second, we use IRT to test
whether the automated scorer is simply a miscalibrated rater: treating it
as a fourth rater alongside the three Delphi judges and $z$-scoring each
criterion across all four (supplementary material), its severity falls
within the range spanned by the three LLM judges---harsher than Claude,
comparable to Gemini---rather than as an outlier. Third, we pool all four
raters' individual ratings into a single graded-response IRT model and
recompute model $\theta$: splitting the pooled fit by rater type
reproduces the same split as before (automated-only $\theta$ favors
HunyuanDiT-3D; Delphi-judge-only $\theta$ favors TRELLIS), and the pooled
estimate narrowly favors TRELLIS ($\theta=0.475$ vs.\ HunyuanDiT-3D's
$0.432$)---but only because the three Delphi judges contribute three times
as many rated instances per model as the single automated score, not
because pooling adjudicates which method is correct. So IRT resolves two
of three candidate explanations (dimension weighting, rater
miscalibration) but not the third: the disagreement is a genuine
measurement-validity question between two instruments, which absent
independent human ground truth we cannot fully settle. That IRT can rule
out two plausible confounds and correctly localize the remaining
disagreement is, we think, a more honest and more useful contribution
than a single adjudicated ranking would have been.

\paragraph{Limitations.}
The Hi3DBench models are all image-to-3D; text-to-3D models (DreamFusion,
Magic3D, etc.)\ may exhibit different psychometric properties.
With only 6 models, the flipped 2PL yields stable prompt-level theta estimates
but model stringency estimates at the extremes (e.g., HunyuanDiT-3D $b=-5.29$)
approach the boundary of reliable estimation; replication with a larger model
cohort is needed.
The DIF analysis is underpowered given the current model distribution (1
non-Objaverse vs.\ 5 Objaverse models).
We also assume, without direct evidence, that stochastic LLM sampling
variance is small relative to genuine disagreement between the three
judge personas; distinguishing these would require repeated ratings from
the same persona at non-zero temperature, which we did not collect, and
is left for future work.
Item-level Delphi--automated-metric agreement is moderate ($r\approx0.53$)---
this is agreement with Hi3DBench's scoring function, not with independent
human labels, which we do not have at item level; future work should
investigate whether longer reasoning chains or domain-specialised judge
personas improve criterion-level calibration, and whether a small human
rating study can adjudicate the HunyuanDiT-3D/TRELLIS disagreement directly.
Hi3DEval includes a second evaluation layer---material scoring---that assesses
physical properties such as albedo, saturation, and metallicness via rendered
video \citep{hi3deval2025}; we analyse only the object-level scores here.
Whether material-level dimensions form a psychometrically distinct factor
from the object-level criteria, and whether material scoring is similarly
susceptible to scale artefacts, are open questions for future work.

%% ─────────────────────────────────────────────────────────────────────────────
\section{Conclusion}
\label{sec:conclusion}

We introduced a Delphi-Psychometric framework for 3D generative model
evaluation, combining structured LLM elicitation with Item Response Theory
and Many-Facet Rasch Modelling.
Applied to Hi3DBench, the two-stream framework produces a coherent set of
findings: automated metric psychometrics reveals that criterion-scale artefacts
dominate benchmark variance (82\% fixed-effect $\eta^2$; 84.7\% under a
random-effects estimate); flipped 2PL IRT recovers stable model
stringency estimates spanning $>8$ logits (HunyuanDiT-3D $b=-5.29$ to
CRM $b=2.83$); 71.6\% of items show $|\Delta p|\geq0.40$ between TripoSR and
the Objaverse-trained models (Objaverse exposure and TripoSR-specific factors
perfectly confounded at $N_\text{non-Objaverse}=1$); and a 30-item
IRT-$\theta$-stratified subset recovers exact model rankings (94\% cost
reduction, LOO-CV $\rho=0.954$).
A cross-level factor analysis extending to Hi3DEval's material-scoring layer
(9 criteria, $N=1{,}730$) finds that object-level and material-level
dimensions are psychometrically distinct constructs with near-zero
cross-loadings (2F CFI=0.883 vs.\ 1F CFI=0.641, $\Delta$BIC$=-1{,}763$),
empirically confirming the hierarchical validity of Hi3DEval's design and
establishing that scores across the two levels should not be aggregated.
Delphi LLM judging cross-validates and complements these findings:
when scale artefacts are removed, model quality emerges as the dominant
variance source (37.3\%); judge reliability is high (ICC$=0.81$--$0.88$);
convergence is strong (68\% inter-judge SD reduction); and Delphi model
rankings agree closely with automated metric rankings ($\rho=0.90$),
providing mutual validation without requiring ground-truth human labels.

\textbf{Relationship to Hi3DEval.}
Hi3DEval \citep{hi3deval2025} establishes \emph{what} to measure---the right
evaluation dimensions and their hierarchical organisation---and solves the
\emph{per-item annotation cost} problem by replacing human labeling with a
video-based automated scorer aligned to human judgment via hybrid
3D-aware representations.
Our framework addresses a complementary efficiency problem: given reliable
per-item scores, \emph{how many items} must be scored to reliably rank
models?
The answer from our IRT analysis is 30, not 510.
Hi3DEval's video-based scoring pipeline (video rendering, InternVideo2-CLIP,
geometric embeddings) need only be run on the 30 IRT-selected prompts rather
than the full pool, reducing total evaluation cost by 94\% without
sacrificing ranking fidelity (LOO-CV $\rho=0.954$).
This composability is especially valuable because Hi3DEval's per-item cost is
substantially higher than static-image metrics.
Beyond efficiency, our framework establishes \emph{how} to measure reliably
and fairly: diagnosing scale artefacts that distort aggregation, placing
models on a calibrated latent scale via IRT, and certifying judge reliability
through structured multi-round elicitation with anonymized disagreement
resolution.
Dimension design and measurement methodology are both necessary for a
benchmark whose outputs can be trusted; the two papers address each in turn.

These findings have direct implications for benchmark design: normalise or
replace criterion-specific scales before aggregation; do not combine
object-level and material-level scores into a single composite; prefer
IRT-$\theta$-stratified or difficulty-stratified item selection; report
per-criterion profiles alongside aggregate scores; and expand the model set
to enable valid DIF testing.
We release all code and the \texttt{tinyEval3D.json} subset file to support
future use.

%% ─────────────────────────────────────────────────────────────────────────────
\bibliographystyle{plainnat}
\bibliography{references}

% Robustness checks are provided in the supplementary material (supplementary.tex).

\end{document}
